\chapter{Proposed System Design}

\section{Overview}
This chapter translates the requirements and analysis artifacts from Chapter~3 into concrete design decisions. It describes the high-level architecture, the internal structure of each service, the API design contracts, the data flow between services, the algorithmic stages within each AI pipeline, the security architecture, and the deployment topology. The design prioritizes modularity, independent scalability of AI workloads, and a clear separation between business logic and inference-heavy components.

\section{Architecture Overview}

\projectname\ follows a \textbf{service-oriented architecture (SOA)} comprising four independently deployable services that communicate over HTTP/REST:

\begin{enumerate}[leftmargin=*]
    \item \textbf{Frontend Service} — A Next.js application that renders the user interface, manages client-side state, handles internationalization (Arabic/English with RTL support), and communicates with the backend and AI services via REST calls.

    \item \textbf{Backend Service} — An Express.js application written in TypeScript that serves as the central orchestrator. It handles authentication, authorization, business logic for e-commerce and community features, proxies requests to AI services, and manages all persistent data through MongoDB.

    \item \textbf{Try-On Service} — A FastAPI application that exposes specialized endpoints for virtual try-on (2D and multiview), body measurement estimation, video preprocessing, and 360-degree reconstruction. It manages connections to external inference servers and local preprocessing pipelines.

    \item \textbf{Chat Service (RAG)} — A FastAPI application that implements the fashion-aware retrieval-augmented generation pipeline. It handles message classification, intent extraction, product retrieval with hybrid search, outfit assembly, and bilingual response generation.
\end{enumerate}

\begin{figure}[H]
    \centering
    \makebox[\textwidth][c]{\includegraphics[width=1.15\textwidth]{figures/arch-overview.png}}
    \caption{High-level service-oriented architecture of \projectname.}
    \label{fig:arch-overview}
\end{figure}

The frontend communicates directly with the backend for all e-commerce, authentication, and community operations. For AI features, two patterns are used: (1)~the backend proxies try-on and measurement requests to the Try-On Service, adding authentication validation and result persistence before forwarding; (2)~the frontend communicates directly with the Chat Service for real-time conversational interactions, as the RAG pipeline maintains its own in-memory conversation history and does not require backend-mediated persistence for the chat flow itself.

\section{Frontend Design}

\subsection{Framework and Routing}
The frontend uses \textbf{Next.js} with the App Router architecture. All user-facing pages are nested under a \texttt{[locale]} dynamic segment that enables locale-based routing. The locale parameter drives language selection (Arabic or English) and layout direction (RTL or LTR). The routing hierarchy is structured as follows:

\begin{verbatim}
app/
  [locale]/
    page.tsx              # Landing / Home
    shop/                 # Product catalog
    product/[id]/         # Product detail
    cart/                 # Shopping cart
    checkout/             # Checkout flow
    orders/               # Order history
    order-success/        # Order confirmation
    profile/              # User profile & settings
    virtual-tryon/        # Try-on interface
    community/            # Posts, designs, feed
    recommendations/      # AI-powered suggestions
    customer-service/     # Support ticket interface
    merchant-owner/       # Merchant dashboard
    admin/                # Admin panel
    login/ | signup/      # Authentication pages
    wishlist/             # Saved items
    dashboard/            # User dashboard
\end{verbatim}

\subsection{State Management and API Communication}
The frontend uses React hooks and context for local state management. API communication is centralized through a typed API client module that handles:
\begin{itemize}[leftmargin=*]
    \item Attaching JWT access tokens from local storage to all authenticated requests.
    \item Intercepting \texttt{401 Unauthorized} responses and checking for a refreshed token in the \texttt{X-New-Access-Token} response header (transparent token rotation).
    \item Providing typed request/response wrappers for each backend endpoint.
\end{itemize}

\subsection{Internationalization (i18n)}
The application supports Arabic and English through Next.js internationalized routing. Each locale has its own translation files loaded at build time. The layout component reads the current locale parameter and applies the appropriate \texttt{dir} attribute (\texttt{rtl} for Arabic, \texttt{ltr} for English) to the root HTML element. All UI text, labels, and messages reference translation keys rather than hardcoded strings.

\section{Backend Design}

\subsection{Module Architecture}
The backend follows a \textbf{layered module pattern} where each domain feature is encapsulated in its own directory under \texttt{src/modules/}. Every module contains four standard files:

\begin{itemize}[leftmargin=*]
    \item \textbf{Routes} (\texttt{*.routes.ts}) — Defines Express router endpoints, attaches middleware (authentication, validation, rate limiting), and maps HTTP methods to controller functions.
    \item \textbf{Controller} (\texttt{*.controller.ts}) — Parses request parameters and body, delegates to the service layer, and formats HTTP responses.
    \item \textbf{Service} (\texttt{*.service.ts}) — Contains business logic, database queries, and integration calls. Services are stateless functions.
    \item \textbf{Model} (\texttt{*.model.ts}) — Defines the Mongoose schema, TypeScript interfaces, indexes, and virtual fields for the module's data entity.
\end{itemize}

An optional \textbf{Schema} file (\texttt{*.schema.ts}) defines Zod validation schemas that are enforced by a generic validation middleware before the request reaches the controller.

\subsection{Middleware Pipeline}
Incoming requests pass through a layered middleware pipeline before reaching route handlers:

\begin{enumerate}[leftmargin=*]
    \item \textbf{Security Headers} — Helmet middleware applies HTTP security headers (Content-Security-Policy, X-Frame-Options, etc.).
    \item \textbf{CORS} — Configured to allow requests only from authorized frontend origins with credentials support.
    \item \textbf{Rate Limiting} — Two tiers: a general API rate limiter and a stricter authentication-specific rate limiter for login, signup, and password reset endpoints.
    \item \textbf{Body Parsing} — JSON and URL-encoded body parsing with configurable size limits.
    \item \textbf{Authentication} — JWT verification with automatic token refresh. The middleware extracts the Bearer token, verifies it against the JWT secret, and attaches the decoded user object to the request. If the access token is expired but a valid refresh token cookie is present, the middleware silently issues a new access token via the \texttt{X-New-Access-Token} response header.
    \item \textbf{Role Authorization} — A configurable role guard that checks whether the authenticated user's role is in the set of permitted roles for the endpoint.
    \item \textbf{Validation} — Zod schema validation middleware that validates request body, query, and params against the module's schema definitions, returning structured error responses on failure.
    \item \textbf{Profanity Guard} — Content moderation middleware applied to user-generated content endpoints.
\end{enumerate}

An \textbf{optional authentication middleware} variant is also provided for endpoints that enhance the experience for authenticated users but remain accessible to anonymous visitors (e.g., product browsing with personalized recommendations).

\subsection{User Caching}
To reduce database load on the authentication hot path, the backend implements an in-memory user cache. On successful JWT verification, the user document is stored in a TTL-bounded LRU cache keyed by user ID. Subsequent requests from the same user within the TTL window are served from cache, bypassing the MongoDB query. Cache entries are invalidated when the user's token version changes (e.g., after password change or forced logout).

\subsection{Error Handling}
A global error handler middleware catches all unhandled exceptions and returns standardized JSON error responses with appropriate HTTP status codes. Domain-specific errors (e.g., \texttt{ServiceUnavailableError}, \texttt{ExternalAPIError}, \texttt{ProcessingError}) are defined as custom exception classes that carry structured metadata for debugging.

\section{Try-On Service Design}

\subsection{Service Architecture}
The Try-On Service is a \textbf{FastAPI} application organized into routers, service clients, preprocessing modules, and a centralized dependency injection container.

\textbf{Routers} define the API surface:
\begin{itemize}[leftmargin=*]
    \item \texttt{/api/tryon} — Single-view 2D virtual try-on using Prova-VTON via a remote inference server.
    \item \texttt{/api/multiangle} — Multi-angle try-on using a multimodal generative model for front, right, left, and back viewpoints.
    \item \texttt{/api/prova-vton} — Multiview pipeline with preprocessing, dataset construction, remote inference, and result retrieval.
    \item \texttt{/api/measurements} — Body measurement estimation using the SHAPY client with size recommendation mapping.
    \item \texttt{/api/video360} — 360-degree video generation from multiview try-on outputs.
    \item \texttt{/api/health} — Aggregated health check across all service clients.
\end{itemize}

\textbf{Service Clients} encapsulate communication with external inference services. Each client follows a consistent pattern:
\begin{itemize}[leftmargin=*]
    \item Configurable base URL, timeouts, and connection pool limits via environment variables.
    \item Automatic retry with exponential backoff for transient failures (HTTP 429, 5xx).
    \item Health check methods for proactive monitoring.
    \item Structured exception hierarchy (\texttt{ServiceUnavailableError}, \texttt{ExternalAPIError}).
\end{itemize}

A \textbf{ClientState} singleton manages the lifecycle of all service clients. It is initialized once at application startup and provides dependency injection to routers via FastAPI's dependency system.

\subsection{Prova-VTON Model Training Architecture}
\label{sec:prova-vton-training-architecture}

The core of the try-on capabilities is powered by Prova-VTON, a custom diffusion-based generative model trained to synthesize realistic garment transfer under pose, appearance, and camera-view constraints. Figure~\ref{fig:vton-model-arch} shows the overall VTON training pipeline. The architecture receives six principal inputs: the target person image \(I_{target}\), the target pose map \(I_{pose}\), an agnostic person image \(I_{agnostic}\), a reference front garment image \(I_{ref\_front}\), a reference back garment image \(I_{ref\_back}\), and camera parameters such as rotation, translation, and intrinsic calibration values. These inputs jointly define what the model should reconstruct during training, which body geometry should be preserved, which visible clothing regions should be replaced, which garment appearance should be transferred, and from which viewpoint the final result should be generated.

\begin{figure}[H]
    \centering
    \makebox[\textwidth][c]{\includegraphics[width=1.15\textwidth]{figures/vton-model-arch.png}}
    \caption{Diffusion-based Prova-VTON training architecture with latent encoding, reference conditioning, pose guidance, camera conditioning, and denoising supervision.}
    \label{fig:vton-model-arch}
\end{figure}

\subsubsection*{Overall Pipeline Overview}
During training, the target image is treated as the reconstruction objective, while the agnostic image, pose map, reference garments, and camera parameters act as conditioning signals. The target image provides the ground-truth visual distribution that the diffusion model must learn to recover after noise is applied in latent space. The agnostic image removes or masks the original garment region from the person, allowing the model to preserve identity, body shape, hair, and background while avoiding direct copying of the original clothing. The pose map contributes structural information about body configuration, and the two reference garment images provide complementary appearance evidence for both front-facing and back-facing garment details. Camera parameters are included because multiview try-on requires the model to understand viewpoint-dependent appearance rather than generating each angle as an unrelated 2D image.

The data flow begins by encoding the target and agnostic images into latent representations. Noise is injected into the target latent at a randomly sampled diffusion timestep. In parallel, the pose map is processed by a PoseGuider, the garment references are encoded by CLIP and ReferenceNet modules, and the camera parameters are projected into a dense camera embedding. The denoising UNet then receives the noisy target latent concatenated with the agnostic latent, resulting in an 8-channel latent input. The conditioning features are injected into the denoising process through attention and feature-control mechanisms. The network is trained to predict the exact Gaussian noise that was added to the clean latent, which allows the reverse diffusion process to reconstruct a clean try-on latent during inference.

\subsubsection*{Latent Space Encoding}
Prova-VTON follows the latent diffusion paradigm. Instead of applying diffusion directly to high-resolution RGB pixels, a variational autoencoder (VAE) encoder \(E(\cdot)\) maps images into a compressed latent space. For the target image, the latent representation is defined as:

\begin{equation}
z = E(I_{target})
\end{equation}

where \(I_{target}\) is the ground-truth target image and \(z\) is its compact latent representation. The agnostic image is encoded in the same latent space to obtain \(z_{agnostic} = E(I_{agnostic})\). The VAE decoder \(D(\cdot)\) is used after denoising to transform the generated latent back into the image domain. Latent diffusion is used because it substantially reduces memory consumption and computational cost compared with pixel-space diffusion. This is particularly important for virtual try-on, where high-resolution garment textures, human body structure, and multiview consistency must be modeled simultaneously. The latent space also provides a more semantically organized representation, enabling the denoising network to focus on garment placement and appearance synthesis rather than low-level pixel reconstruction alone.

\subsubsection*{Reference Feature Extraction}
The reference branch is responsible for preserving garment identity. A CLIP-based reference encoder converts \(I_{ref\_front}\) and \(I_{ref\_back}\) into embedding vectors that capture high-level garment semantics, including category, silhouette, color, and texture cues. These embeddings are then further processed by ReferenceNet, an appearance encoder that produces spatial reference features suitable for injection into the denoising UNet. The use of both front and back garment images is important for multiview training because a single product image cannot reliably represent garment regions that become visible from side or rear viewpoints.

Camera conditioning is computed by passing the camera parameters through an embedding module, implemented conceptually as a Fourier or multilayer-perceptron projection. The resulting vector \(c_{cam}\) allows the network to associate visual evidence with the requested viewpoint. This design reduces inconsistencies such as front-view textures appearing incorrectly in back-view outputs, because the denoiser receives an explicit representation of the camera pose.

\subsubsection*{Forward Diffusion Process}
Training uses a denoising diffusion probabilistic model (DDPM) scheduler. At each iteration, a timestep \(t\) is randomly sampled from the training diffusion horizon, and Gaussian noise \(\epsilon\) is sampled from a standard normal distribution:

\begin{equation}
\epsilon \sim \mathcal{N}(0, I)
\end{equation}

The scheduler then produces a noisy latent \(z_t\) according to:

\begin{equation}
z_t = \sqrt{\bar{\alpha}_t} z + \sqrt{1 - \bar{\alpha}_t}\epsilon
\end{equation}

In this equation, \(z\) is the clean latent representation of the target image, \(z_t\) is the noised latent at timestep \(t\), \(\epsilon\) is the Gaussian noise added to the latent, and \(\bar{\alpha}_t\) is the cumulative product of the scheduler's noise-retention coefficients up to timestep \(t\). Small values of \(t\) preserve much of the original signal, while larger values of \(t\) make the latent increasingly noisy. Random timestep sampling teaches the model to denoise from different corruption levels, which is necessary for the iterative reverse diffusion process used during inference.

\subsubsection*{Conditioning Modules}
The conditioning modules provide the denoising network with information that cannot be inferred from the noisy latent alone. The PoseGuider converts \(I_{pose}\) into pose features that encode body layout, limb orientation, and spatial correspondence between the person and the intended garment region. These generated pose features guide the model toward anatomically plausible garment placement and reduce artifacts around arms, torso boundaries, and self-occlusions.

ReferenceNet produces generated reference features from the front and back garment images. These features carry the garment's appearance and are responsible for preserving details such as fabric color, collar structure, sleeves, seams, logos, and texture patterns. The camera embedding provides viewpoint awareness, enabling the same garment to be synthesized consistently across different camera directions. Together, these conditioning signals separate the main factors of variation: pose defines where the garment should appear, reference features define what the garment should look like, and camera features define how it should be observed.

\subsubsection*{UNet Denoising Network}
The denoising backbone is a conditional UNet with an encoder-decoder structure, skip connections, and cross-attention blocks. Its input is formed by concatenating the noisy target latent \(z_t\) and the agnostic latent \(z_{agnostic}\) along the channel dimension. Since each latent contains four channels, the concatenated tensor forms an 8-channel input. This input design allows the model to compare the corrupted target representation with the preserved person context from the agnostic image. The encoder progressively extracts abstract features, the bottleneck integrates global conditioning, and the decoder reconstructs a noise prediction at the original latent resolution. Skip connections preserve spatial detail across scales, which is essential for maintaining face, hands, body boundaries, and garment edges.

Pose features, reference features, and camera embeddings are injected through attention blocks. In the cross-attention formulation, the UNet hidden states provide the queries, while conditioning features provide keys and values:

\begin{equation}
\operatorname{Attention}(Q,K,V) = \operatorname{Softmax}\left(\frac{QK^T}{\sqrt{d}}\right)V
\end{equation}

Here, \(Q\) represents query vectors derived from the current UNet feature maps, \(K\) represents key vectors derived from conditioning features, \(V\) represents value vectors containing the conditioning information to be integrated, and \(d\) is the key dimensionality used for normalization. In Prova-VTON, this means that each spatial region of the denoising feature map can selectively attend to pose structure, garment appearance, and camera context. For example, pixels corresponding to the upper torso can attend strongly to jacket texture features, while arm regions can attend more heavily to pose features to maintain correct sleeve alignment.

\subsubsection*{Reference Attention Mechanism}
The reference attention mechanism is implemented through a Reference Control Writer and a Reference Control Reader. During the reference pass, ReferenceNet extracts garment-aware features from the front and back reference images, and the writer stores these features in a control buffer associated with the attention layers. During the main denoising pass, the reader retrieves the stored features and supplies them to the UNet attention blocks. This separation allows the model to encode garment appearance once and reuse it throughout denoising, rather than forcing the UNet to infer garment identity only from a single global embedding.

This mechanism is important for preserving garment appearance. Direct concatenation of reference images would require the network to learn alignment and feature selection implicitly, which is difficult when garment and body poses differ. Reference attention instead allows the denoiser to retrieve appearance cues dynamically at multiple resolutions. As a result, fine-grained garment details can be maintained while still adapting the garment to the target pose and camera view.

\subsubsection*{Noise Prediction and Training Objective}
The UNet is trained to predict the injected noise:

\begin{equation}
\hat{\epsilon} = \operatorname{UNet}(z_t, z_{agnostic}, F_{pose}, F_{ref}, c_{cam}, t)
\end{equation}

where \(\hat{\epsilon}\) is the predicted noise, \(F_{pose}\) denotes pose features, \(F_{ref}\) denotes reference garment features, \(c_{cam}\) denotes the camera embedding, and \(t\) denotes the diffusion timestep. The training loss is the mean squared error between the true noise and the predicted noise:

\begin{equation}
L_{MSE} = \mathbb{E}\left[\lVert \epsilon - \hat{\epsilon} \rVert^2\right]
\end{equation}

By minimizing this objective, the model learns how to remove noise from latent representations while respecting the agnostic body context and all conditioning signals. Since the clean latent \(z\) corresponds to the target try-on image, accurate noise prediction implicitly teaches the network to reconstruct clean latent representations that contain the correct person identity, garment appearance, pose alignment, and camera viewpoint.

\subsubsection*{Training and Inference Separation}
The architecture distinguishes clearly between training and inference. During training, the target image is available and is used to produce \(z\), sample \(z_t\), and compute the supervised noise-prediction loss. The ground-truth noise \(\epsilon\) is known because it is sampled by the training algorithm. During inference, the target image and ground-truth noise are not available. Instead, the process begins from random Gaussian latent noise, and the trained UNet iteratively predicts and removes noise while conditioned on the user's agnostic image, pose map, garment references, and camera parameters. After the final reverse diffusion step, the VAE decoder maps the denoised latent into the final try-on image:

\begin{equation}
\hat{I}_{tryon} = D(\hat{z}_0)
\end{equation}

where \(\hat{z}_0\) is the final denoised latent and \(\hat{I}_{tryon}\) is the synthesized try-on output. This reconstruction stage converts the learned latent representation into a human-viewable image suitable for display in the e-commerce interface or for later use in multiview and 360-degree preview workflows.

\subsection{Multi-Angle Try-On Pipeline Design}
The multi-angle try-on client implements a two-stage pipeline for producing realistic virtual try-on results across multiple viewpoints:

\textbf{Stage 1 --- Garment Preparation.} Before try-on synthesis, the system determines whether the garment reference image contains a model wearing the garment or is a flat product shot. The decision is made through a three-step heuristic:
\begin{enumerate}[leftmargin=*]
    \item Check for transparent background (alpha channel analysis) --- transparent images are classified as flat product shots.
    \item Detect skin-tone pixels using RGB-space rules and check for portrait aspect ratio --- images with significant skin content in portrait orientation are classified as model-worn.
    \item If the source mode is ambiguous (\texttt{auto}), the heuristic result is used; explicit modes (\texttt{flat\_product}, \texttt{model\_worn}) override the heuristic.
\end{enumerate}

For model-worn garment references, an extraction step uses the multimodal model to isolate the garment from the person, producing a clean product image before try-on synthesis.

\textbf{Stage 2 --- Try-On Synthesis.} The prepared garment image and the person image are submitted to the generative model with a composed prompt. The prompt is dynamically constructed from a base template augmented with control parameters:
\begin{itemize}[leftmargin=*]
    \item \textbf{Garment type} (upper/lower/dress/shoes) --- instructs the model on which body region to modify.
    \item \textbf{Shirt placement} (tucked/untucked) --- controls hem behavior.
    \item \textbf{View hint} (front/right/left/back) --- preserves the camera orientation of the person image.
\end{itemize}

For multi-angle try-on, the system processes multiple viewpoints concurrently using an asyncio semaphore (concurrency limit of 2) to balance throughput against API rate limits. Garment preparation is performed once and shared across all views.

\subsection{Multiview (Prova-VTON) Pipeline Design}
The multiview pipeline is designed for generating view-consistent try-on results across all angles. It operates in five stages:

\begin{enumerate}[leftmargin=*]
    \item \textbf{Preprocessing} — Input video or images are processed through YOLOv8 for person detection, Detectron2 for DensePose map generation, and a human parsing network for segmentation masks. Results are organized into a structured directory hierarchy.

    \item \textbf{Cloth Preparation} — Garment front and back images are saved, resized, and organized into the dataset structure expected by the inference model.

    \item \textbf{Dataset Construction} — Person preprocessing outputs and cloth images are assembled into a ZIP archive following the model's expected directory format, ready for upload.

    \item \textbf{Remote Inference} — The ZIP is uploaded to the inference server. The service polls for job completion with configurable timeouts.

    \item \textbf{Result Retrieval} — Completed results are downloaded, extracted, and made available through the API. The multiview outputs can optionally be fed into the Splatfacto 3D Gaussian splatting pipeline for 360-degree interactive reconstruction.
\end{enumerate}

\begin{figure}[H]
\centering
\begin{tikzpicture}[
    node distance=1.5cm and 1cm,
    process/.style={rectangle, minimum width=3.5cm, minimum height=1cm, text centered, draw=provablue, fill=provalight, font=\small\bfseries, rounded corners},
    arrow/.style={thick,->,>=stealth, provablue}
]

% Nodes
\node (n1) [process] {1. Preprocessing (YOLO, DensePose, SCHP)};
\node (n2) [process, below=of n1] {2. Cloth Preparation (RGB/JPEG Format)};
\node (n3) [process, below=of n2] {3. Dataset Construction (ZIP Archive)};
\node (n4) [process, below=of n3] {4. Remote Inference (GPU Cluster)};
\node (n5) [process, below=of n4] {5. Result Retrieval \& 3D Splatting};

% Arrows
\draw [arrow] (n1) -- (n2);
\draw [arrow] (n2) -- (n3);
\draw [arrow] (n3) -- (n4);
\draw [arrow] (n4) -- (n5);

\end{tikzpicture}
\caption{The five-stage architecture of the Prova-VTON multiview pipeline.}
\label{fig:prova-vton-pipeline}
\end{figure}

\subsection{Measurement Pipeline Design}
The measurement pipeline transforms a single front-facing photograph into body measurements and size recommendations through the following stages:

\begin{enumerate}[leftmargin=*]
    \item \textbf{Input Validation} — The router validates image format, file size, height (in cm), weight (in kg), and gender parameters.

    \item \textbf{SHAPY Inference} — The image and metadata are sent to the SHAPY service, which performs SMPL-X body mesh regression and returns raw shape parameters and circumference estimates.

    \item \textbf{Weight-Based Correction} — If the predicted body mass deviates significantly from the user-provided weight, the router applies a proportional correction factor to circumference measurements. This compensates for the inherent ambiguity of estimating depth information from a single 2D image.

    \item \textbf{Derived Measurements} — Additional body dimensions (shoulder width, arm length, trouser length, thigh circumference) are computed from the corrected circumferences using ISO~7250 anthropometric ratios.

    \item \textbf{Size Recommendation} — Chest and waist circumferences are mapped to regional size charts (EU, US, UK, IT, FR, DE) with configurable fit preferences (regular, slim, oversized). The fit preference applies a scaling factor: slim reduces the target circumference, oversized increases it.

    \item \textbf{Confidence Scoring} — Each measurement receives a confidence score based on the quality of the input image, the magnitude of corrections applied, and the consistency between predicted and reported anthropometric data.
\end{enumerate}

\section{Chat Service (RAG) Design}

\subsection{Fashion Recommendation RAG Architecture}
\label{sec:rag-architecture}

The Chat Service implements a fashion recommendation architecture based on retrieval-augmented generation. Its purpose is to transform a user's natural language shopping request into grounded product and outfit recommendations drawn from the actual catalog. Figure~\ref{fig:rag-model-arch} shows the complete RAG recommendation workflow. The design separates retrieval from generation: retrieval identifies products that satisfy the user's constraints, while generation converts those retrieved products into a natural, bilingual explanation. This separation is essential in an e-commerce setting because the assistant must not merely produce plausible fashion advice; it must recommend items that exist in inventory and satisfy availability, category, price, style, and user-preference constraints.

\begin{figure}[H]
    \centering
    \makebox[\textwidth][c]{\includegraphics[width=1.15\textwidth]{figures/rag-model-arch.png}}
    \caption{Fashion recommendation RAG architecture from user message classification to hybrid retrieval, outfit scoring, bilingual response composition, and final output.}
    \label{fig:rag-model-arch}
\end{figure}

\subsubsection*{User Query Processing}
The workflow begins when the user submits a message through the AI fashion assistant. The Message Classifier determines whether the message is a product-search request, a greeting, a follow-up, a clarification, or a non-fashion query. This early decision prevents unnecessary retrieval for messages that do not require catalog access. For example, greetings and general conversational turns are routed to a Quick Response module, which returns a concise response without invoking the retrieval and ranking pipeline. Non-fashion messages are also handled through this fast path, keeping the assistant focused on the shopping domain.

When the classifier detects a product-search intent, the message is forwarded to the intent extraction stage. This creates a clear distinction between conversational inference and recommendation inference. Conversational inference handles lightweight replies, while recommendation inference performs structured interpretation, database retrieval, semantic ranking, outfit assembly, and response generation.

\subsubsection*{Intent Extraction}
The Intent Extractor converts natural language into a structured fashion intent. The extracted intent may include style, occasion, season, gender, garment type, colors, size, budget, brand preferences, fit, material, and additional attributes. For example, the message ``I need a modest summer outfit for a casual university day'' can be converted into an intent containing \texttt{style=modest}, \texttt{season=summer}, \texttt{occasion=casual/university}, \texttt{gender} inferred from the user profile when available, and relevant garment categories such as tops, trousers, and light outerwear. Similarly, ``Find me a black formal dress under 2000 EGP'' maps to a structured query with color, occasion, garment type, and price constraints.

This structured representation is necessary because product catalogs are not searched effectively by raw natural language alone. Hard constraints such as price range, size, gender, and stock availability must be represented explicitly so that retrieval does not return unsuitable products. At the same time, softer preferences such as ``elegant,'' ``streetwear,'' or ``minimal'' are preserved for semantic ranking because they may not correspond to exact database fields.

\subsubsection*{Hybrid Retrieval Layer}
The Hybrid Retriever combines MongoDB metadata filtering with semantic retrieval using Sentence-BERT. The first mechanism applies deterministic filters over catalog fields. MongoDB queries restrict the candidate set according to attributes such as category, gender, color, size availability, price range, brand, and stock status. This stage improves precision and computational efficiency by eliminating products that violate explicit user constraints before semantic ranking is performed.

The second mechanism ranks the remaining candidates using dense embeddings. The user's query is encoded as:

\begin{equation}
q = \operatorname{SBERT}(\operatorname{Query})
\end{equation}

Each product is represented by an embedding computed from its searchable textual fields:

\begin{equation}
p_i = \operatorname{SBERT}(\operatorname{Product}_i)
\end{equation}

The semantic relevance between the user query and product \(i\) is measured using cosine similarity:

\begin{equation}
\operatorname{Similarity}(q,p_i) = \frac{q \cdot p_i}{\lVert q \rVert \lVert p_i \rVert}
\end{equation}

Here, \(q\) is the query embedding, \(p_i\) is the embedding of the \(i\)-th product, \(q \cdot p_i\) is their dot product, and \(\lVert \cdot \rVert\) denotes vector magnitude. Cosine similarity is appropriate because it measures directional closeness in embedding space, allowing semantically related descriptions to match even when they do not share exact keywords. The hybrid design improves recommendation quality by combining the precision of structured filtering with the flexibility of semantic understanding.

\subsubsection*{Retrieval Phase}
The retrieval phase generates candidate products, ranks them, and selects the top products for outfit construction. Candidate generation begins with metadata filtering so that the system respects explicit constraints. Dense ranking then orders candidates according to semantic relevance. Products with low similarity can be discarded, while the highest-ranked products are passed forward as top candidate products. This staged design reduces hallucination risk because the response composer receives only verified catalog items rather than unconstrained model-generated suggestions.

\subsubsection*{Outfit Assembler}
The Outfit Assembler builds complete outfits from the retrieved items. Rather than recommending isolated products independently, the assembler groups items according to category compatibility and styling logic. For example, a top can be paired with trousers or a skirt, outerwear can be added when appropriate for the season, and shoes or accessories can be included when available. The assembler also avoids incompatible combinations, such as pairing two primary tops in the same outfit unless one is classified as an outer layer. This step translates product retrieval into a more useful fashion recommendation because users typically seek complete looks rather than disconnected catalog entries.

\subsubsection*{Outfit Scoring Engine}
After candidate outfits are assembled, a scoring engine evaluates them using a six-dimensional framework. The overall outfit score can be expressed as:

\begin{equation}
\operatorname{Score}(\operatorname{outfit}) = \sum_i w_i s_i
\end{equation}

In this equation, \(s_i\) is the score assigned to the outfit on dimension \(i\), and \(w_i\) is the weight expressing the importance of that dimension. Suitable dimensions include style compatibility, color harmony, occasion suitability, seasonal suitability, trend alignment, and user preference matching. Style compatibility measures whether the selected products belong to a coherent aesthetic. Color harmony checks whether the palette is visually consistent. Occasion suitability evaluates whether the outfit matches the intended use case, such as formal events, university, work, or casual wear. Seasonal suitability considers fabric weight, sleeve length, and layering. Trend alignment reflects current or catalog-defined fashion signals, while user preference matching incorporates profile information, previous conversation context, and explicit constraints.

The weighted formulation allows the system to adapt recommendations to different situations. For a wedding query, occasion suitability may receive a high weight; for a user asking for daily basics, comfort, color neutrality, and preference matching may become more important. The scoring engine therefore provides a principled way to rank outfits beyond simple product-level relevance.

\subsubsection*{Outfit Grouping and Ranking}
The grouped outfits are ranked according to their total scores and diversity constraints. Grouping similar outfits prevents the assistant from returning several near-identical combinations that differ only by minor product substitutions. Ranking then selects the best recommendations according to the scoring function, while diversity ensures that the final response contains meaningful alternatives. This step improves user experience by presenting a small set of strong, distinct recommendations rather than a long list of products.

\subsubsection*{Bilingual Response Composer}
The Bilingual Response Composer receives the selected outfits, product metadata, the original user query, conversation history, and available user profile information. It generates a human-readable response in English or Arabic, depending on the user's language preference or detected message language. The response is grounded in the retrieved products and includes explanations for why each outfit was selected. For instance, it may explain that a linen shirt was recommended because it matches a summer requirement, or that a darker trouser balances a bright top for a more formal appearance.

This generation stage is intentionally placed after retrieval and scoring. The language model is not responsible for inventing products; it is responsible for communicating catalog-grounded recommendations clearly. This design supports bilingual interaction while preserving factual alignment with the product database.

\subsubsection*{Final System Output}
The final output returned to the user consists of one or more recommended outfits, associated product references, and explanatory text. End-to-end, information moves from an unstructured user message to a classified intent, then to a structured fashion query, through hybrid retrieval, outfit assembly, scoring, grouping, and finally bilingual response generation. If the message is not a product-search request, the system bypasses the retrieval path and returns a quick response. If it is a product-search request, the full RAG workflow ensures that recommendations are relevant, available, stylistically coherent, and explainable.

\subsection{Conversation History}
The service maintains per-user conversation history in an in-memory LRU (Least Recently Used) store. The store is keyed by \texttt{user\_id} (primary) or \texttt{session\_id} (fallback), with configurable limits: up to 5{,}000 distinct keys and 20 turns (user + assistant messages) per key. Older entries are evicted when limits are reached. This enables multi-turn conversations where the assistant can reference previous queries without requiring external storage.

\subsection{Knowledge Base Integration}
A fashion-specific knowledge base provides profile-aware styling guidance. The knowledge base contains rules indexed by body type, height range, and gender that map to recommended silhouettes, colors, and fit preferences. When the user's profile includes height, weight, and gender, the retriever consults the knowledge base to augment recommendations with personalized styling advice.

\subsection{Embedding Strategy}
Product embeddings are pre-computed using Sentence-BERT (\texttt{all-MiniLM-L6-v2}) over a concatenation of product name, description, brand, tags, styles, and occasion fields. Embeddings are stored alongside product documents in MongoDB. A dedicated admin endpoint (\texttt{POST /precompute-embeddings}) triggers bulk re-computation when the product catalog is updated.

\section{API Design}

The system exposes RESTful APIs following consistent design conventions:

\subsection{Backend API Conventions}
\begin{itemize}[leftmargin=*]
    \item \textbf{URL Structure}: Resources follow a noun-based hierarchy (e.g., \texttt{/api/products}, \texttt{/api/orders/:id}, \texttt{/api/community/posts}).
    \item \textbf{Authentication}: Protected endpoints require a \texttt{Bearer} token in the \texttt{Authorization} header. Public endpoints use the optional auth middleware.
    \item \textbf{Response Format}: All responses follow a consistent JSON structure with \texttt{success} (boolean), \texttt{msg} (human-readable message), and \texttt{data} (payload) fields.
    \item \textbf{Error Format}: Errors include \texttt{success: false}, a machine-readable \texttt{code} (e.g., \texttt{TOKEN\_EXPIRED}), and a descriptive \texttt{msg}.
    \item \textbf{Pagination}: List endpoints support \texttt{page} and \texttt{limit} query parameters with default values.
    \item \textbf{Validation}: Request bodies are validated against Zod schemas before reaching controllers.
\end{itemize}

\subsection{Try-On Service API}
The Try-On Service accepts multipart form data for image uploads. Key endpoints include:

\begin{longtable}{l c p{8cm}}
\toprule
\textbf{Endpoint} & \textbf{Method} & \textbf{Description} \\
\midrule
\endfirsthead
\toprule
\textbf{Endpoint} & \textbf{Method} & \textbf{Description} \\
\midrule
\endhead
\bottomrule
\endfoot

\texttt{/api/tryon/try-on} & POST & Submit person and garment images for 2D try-on using Prova-VTON. \\
\texttt{/api/multiangle/tryon} & POST & Single-view try-on with garment type, source mode, and placement controls. \\
\texttt{/api/multiangle/multi} & POST & Multi-angle try-on across front, right, left, and back viewpoints. \\
\texttt{/api/measurements/measure} & POST & Body measurement estimation from a front-facing photograph. \\
\texttt{/api/prova-vton/preprocess} & POST & Preprocess person video/images for the multiview pipeline. \\
\texttt{/api/prova-vton/upload-cloth} & POST & Upload garment front and back images. \\
\texttt{/api/prova-vton/build-dataset} & POST & Assemble preprocessed data into inference-ready format. \\
\texttt{/api/prova-vton/infer} & POST & Submit dataset for multiview inference. \\
\texttt{/api/prova-vton/results} & GET & Retrieve completed inference results. \\
\texttt{/api/health} & GET & Aggregated health status across all service clients. \\
\end{longtable}

\subsection{Chat Service API}
\begin{longtable}{l c p{8cm}}
\toprule
\textbf{Endpoint} & \textbf{Method} & \textbf{Description} \\
\midrule
\endfirsthead
\toprule
\textbf{Endpoint} & \textbf{Method} & \textbf{Description} \\
\midrule
\endhead
\bottomrule
\endfoot

\texttt{/chat} & POST & Send a natural language message through the RAG pipeline. Accepts \texttt{message}, \texttt{userId}, \texttt{sessionId}, and \texttt{lang}. \\
\texttt{/health} & GET & Service health check with uptime reporting. \\
\texttt{/precompute-embeddings} & POST & Admin utility to re-compute all product embeddings. \\
\end{longtable}

\section{Security and Privacy Design}

\subsection{Authentication Flow}
The system implements a dual-token authentication strategy:
\begin{itemize}[leftmargin=*]
    \item \textbf{Access Token} — A short-lived JWT (configurable TTL, typically 15 minutes) containing \texttt{userId} and \texttt{tokenVersion}. Sent in the \texttt{Authorization: Bearer} header.
    \item \textbf{Refresh Token} — A longer-lived JWT stored as an HTTP-only cookie. Used to silently rotate access tokens when they expire.
    \item \textbf{Token Versioning} — Each user has a \texttt{tokenVersion} counter. Changing the password or performing a forced logout increments this version, immediately invalidating all existing tokens.
\end{itemize}

\subsection{Password Security}
Passwords are hashed using \textbf{bcrypt} with a configurable salt rounds parameter before storage. Raw passwords are never persisted or logged. The password field is excluded from all database queries by default using Mongoose's \texttt{select: false} option.

\subsection{Rate Limiting}
Two tiers of rate limiting protect the system:
\begin{itemize}[leftmargin=*]
    \item \textbf{General API rate limiter} — Applied to all endpoints with a generous window to prevent abuse while allowing normal usage patterns.
    \item \textbf{Authentication rate limiter} — Stricter limits on login, signup, and password reset endpoints to mitigate brute-force attacks. A separate \textbf{password reset rate limiter} is applied to the forgot-password, verify-code, and reset-password flow.
\end{itemize}

\subsection{Input Validation and Sanitization}
All user inputs are validated at two levels:
\begin{itemize}[leftmargin=*]
    \item \textbf{Schema validation} (Zod) — Ensures type safety, required fields, string length limits, and format constraints (e.g., email format, phone number patterns).
    \item \textbf{Content moderation} — A profanity guard middleware scans user-generated text content before it reaches the database.
\end{itemize}

\subsection{Data Privacy}
Uploaded images for try-on and body measurements are processed in-memory where possible and stored in Cloudinary with access-controlled URLs. Try-on results are associated with authenticated user accounts, ensuring that results are only accessible to the user who generated them. The body measurement pipeline does not retain raw images beyond the inference session unless explicitly saved by the user.

\section{Deployment Architecture}

The system is designed for deployment across multiple environments:

\begin{itemize}[leftmargin=*]
    \item \textbf{Frontend} — Deployed as a containerized Next.js application behind a reverse proxy (Nginx or Vercel). Static assets are served from a CDN.

    \item \textbf{Backend} — Runs as a Node.js process with PM2 for process management and automatic restarts. Connects to a managed MongoDB Atlas cluster.

    \item \textbf{Try-On Service} — Runs as a Uvicorn ASGI server. GPU-intensive inference (Prova-VTON) is offloaded to remote GPU servers accessed via HTTP. The local service handles preprocessing, orchestration, and result delivery.

    \item \textbf{Chat Service} — Runs as a Uvicorn ASGI server. The Sentence-BERT embedding model runs locally on CPU. LLM inference is delegated to an external API endpoint.

    \item \textbf{Database} — MongoDB Atlas provides the primary data store with automated backups, replica sets for high availability, and Atlas Search for text indexing.

    \item \textbf{Media Storage} — Cloudinary serves as the cloud media storage provider for product images, profile photos, try-on results, and community post images.
\end{itemize}

\section{Summary}
This chapter has presented the system design for \projectname, covering the service-oriented architecture, the internal design of each service (frontend routing and i18n, backend module pattern and middleware pipeline, try-on multi-stage pipelines, and RAG pipeline stages), the API contracts, the security architecture, and the deployment topology. These design decisions directly implement the requirements defined in Chapter~3 and provide the blueprint for the implementation described in Chapter~5.
